% !TeX root = ../main.tex
% ---------------------------------------------------------------
\section{Discussion}
\label{sec:discussion}
% ---------------------------------------------------------------

The comparison shows that spectral researcher scores and the Stanford composite measure different things. The Stanford list ranks researchers by the accumulated citation record of a career or of a year, and the people at its top have long publication lists and thousands of citations. The spectral scores rank researchers by the typical standing of their papers, judged either by the influence of the journals and institutions that produced them or by the influence of the sources that cite them. Among researchers already on the Stanford list the two are only weakly correlated, and when both are drawn up as top-2\% lists within subfields they select largely different people. The correlation with Stanford's size-independent citations per paper is higher, which locates most of the difference in the treatment of volume rather than in the data sources.

Neither view is the correct one. A cumulative indicator answers the question of whose work has, in total, been most cited, and it is the natural basis for recognising sustained contributions. A mean answers the question of how the typical paper of a researcher is received, and it does not favour long careers or large outputs. Its weakness is the mirror image of the cumulative indicator's. With a threshold of five works, a short record of papers cited from influential sources can place a researcher in the top 2\%, and the median researcher in our citer-score list has only \NTopMedWorksCit{} works. A score of this kind is better read with its supporting counts, and a list that is to be used for recognition would need a higher volume threshold or a ranking that combines the mean with the number of works.

Weighting citations by the influence of their source changes the role of self-citation. Removing self-citations reorders researchers only slightly, and for researchers with heavy self-citation it raises their citer scores rather than lowering them, because their self-citations tend to appear in less influential venues. A mean of citer influence can be raised by citations only if those citations come from sources more influential than the researcher's existing citers, which is precisely what a self-citing researcher or a citation cartel cannot easily arrange. This is the main practical difference from count-based indicators, whose authors must remove self-citations and still caution that citation farms may remain \citep{ioannidisStandardizedCitationMetrics2019a}. The citer score does not detect manipulation, but it gives manipulation little to gain.

The country composition of the lists deserves caution. The publication score inherits the influence of a researcher's journals and institutions, and these are concentrated in the United States and a few other countries, so a ranking on publication score reproduces the hierarchy of institutions as much as it ranks individuals. The citer score is less tied to the researcher's own institution, and its list is closer to the Stanford list in country composition. In both, researchers in China and India are much less represented than among all eligible researchers. This may reflect lower influence of the sources that cite their work, the publication of much of their work in sources that are not retained in the field rankings, or differences in OpenAlex author disambiguation for common names, and our data cannot separate these explanations.

\subsection{Limitations}

Several limitations bear on the comparison, and we note the likely direction of each. First, the matching of Stanford researchers to OpenAlex is by name, location and volume. It recovers four fifths of the career table, and the hand check suggests that most matches are correct, but matching is harder for common names and for early-career researchers, so the matched set probably over-represents senior researchers with distinctive names. Second, OpenAlex author profiles are algorithmic. Split profiles lower a researcher's work count and may remove them from eligibility, and merged profiles mix the papers of different people, which pulls their scores towards the mean. Third, the corpus excludes papers with 30 or more authors and starts in 2000, which removes much of the citation record of researchers in large collaborations and of senior researchers' early work. Fourth, every work is scored by the influence of sources and institutions in 2020--2024, so a paper of 2005 is judged by its journal's present standing rather than its standing at the time. Finally, a citer score rests on the citing works that carry a score, and citations from sources outside the field rankings, probably more common in the humanities and in national-language publishing, are ignored.

\subsection{Conclusion}

We have computed spectral scores for \NOaEligible{} researchers in OpenAlex who meet the Stanford volume baseline, under full and fractional counting, with and without self-citations, for career and single-year citation windows. The scores are open, field-normalised through the influence of sources and institutions, and size-independent. They agree with the Stanford indicators about self-citation, but they rank researchers differently, because they measure the typical reception of a researcher's papers rather than its total. An open alternative to the Stanford list should therefore not be judged by how closely it reproduces that list. The more useful comparison is between what each list rewards. Further work should combine the mean score with volume in a single ranking, test whether the citer score is less affected than $c$ by citation farms identified independently, and repeat the analysis with the sentinel treatment of citations from outside each field.
